\documentclass[11pt,a4paper]{article}
\usepackage[utf8]{inputenc}
\usepackage[margin=1in]{geometry}
\usepackage{hyperref}
\usepackage{array}
\usepackage{booktabs}
\usepackage{enumitem}
\usepackage{titlesec}
\usepackage{xcolor}

\definecolor{cudenvergold}{HTML}{8A724F}
\definecolor{darkgrey}{HTML}{333333}

\titleformat{\section}{\large\bfseries\color{darkgrey}}{\thesection}{1em}{}[{\color{cudenvergold}\titlerule[1.5pt]}]
\titleformat{\subsection}{\normalsize\bfseries\color{darkgrey}}{\thesubsection}{1em}{}

\hypersetup{
    colorlinks=true,
    linkcolor=darkgrey,
    filecolor=magenta,      
    urlcolor=cudenvergold,
    pdftitle={ISMG 2050 Fall 2026 Syllabus},
    pdfpagemode=FullScreen,
}

\begin{document}

\begin{center}
    {\LARGE \bfseries University of Colorado Denver} \\
    \vspace{2mm}
    {\Large \bfseries Business School} \\
    \vspace{4mm}
    {\color{cudenvergold}\hrule height 2pt}
    \vspace{3mm}
    {\large \bfseries ISMG 2050: Business Problem Solving Tools} \\
    {\large \bfseries Course Syllabus --- Fall 2026} \\
\end{center}

\vspace{0.5cm}

\section*{Course Logistics \& Instructor Information}
\begin{itemize}[leftmargin=*]
    \item \textbf{Professor:} Dr. Mohamad Saleh
    \item \textbf{Email:} \href{mailto:mohamad.saleh@ucdenver.edu}{mohamad.saleh@ucdenver.edu} 
    \item \textbf{Office Location:} Business School Building, Room BUSB-4107 or Zoom
    \item \textbf{Office Hours (Zoom):} Mondays 11:00 AM -- 12:30 PM, or by appointment
    \item \textbf{Office Zoom Link:} \url{https://ucdenver.zoom.us/s/2255916025}
    \item \textbf{Student Success Coach (SSC):} Abdelkarim Mansour
    \item \textbf{Coach Email:} \href{mailto:abdelkarim.mansour@ucdenver.edu}{abdelkarim.mansour@ucdenver.edu}
    \item \textbf{Class Schedule:} Tuesday \& Thursday, 3:30 PM to 4:45 PM
    \item \textbf{Classroom Location:} Business School Building, Room 2501
    \item \textbf{Term:} Fall 2026 (08/25/2026 -- 12/12/2026)
\end{itemize}

\section*{Course Description}
ISMG 2050 focuses on the technology and problem-solving skills necessary for students to succeed both at school and in the business world. This course teaches how to make business decisions using spreadsheets, data analysis, and web tools. Students solve problems in statistics, accounting, finance, marketing, management, and information systems. This hands-on course specifically utilizes Microsoft Excel and Tableau to provide students with problem-solving methods and tools necessary to succeed in the business community.

As a business core course, a grade of a 'C-' or better must be earned to satisfy Business graduation requirements and prerequisites for higher-level business courses. 

\section*{Course Goals \& Learning Objectives}
By the end of this course, you will be able to:
\begin{itemize}
    \item Obtain basic proficiency in Microsoft Excel to provide a foundation for future business courses and professional development.
    \item Write formulas containing functions and operators in spreadsheet cells.
    \item Determine effective data display with formatting features, charts, and pivot tables.
    \item Write formulas for rules used in business decision-making situations.
    \item Understand introductory data management concepts and terminology.
    \item Obtain basic proficiency in Tableau to develop business intelligence capabilities.
    \item Transform raw data into meaningful visualizations and insights.
    \item Create interactive dashboards and stories, and handle multiple data sources in Tableau.
\end{itemize}

\section*{Grading Criteria \& Weighting}
The finalized course grade is calculated based on the following weights:
\begin{center}
    \begin{tabular}{lc}
        \toprule
        \textbf{Item} & \textbf{Weight} \\
        \midrule
        Homework Assignments & 70\% \\
        Exams (Mid-Term \& Final Exam) & 30\% \\
        \bottomrule
    \end{tabular}
\end{center}

\subsection*{Undergraduate Letter Grade Distribution}
\begin{center}
    \begin{tabular}{llp{7cm}}
        \toprule
        \textbf{Letter Grade} & \textbf{GPA Value} & \textbf{Percentage Range} \\
        \midrule
        A  & (4.0) & 93\% -- 100\% (Superior/Excellent) \\
        A- & (3.7) & 90\% -- 92.999\% \\
        B+ & (3.3) & 87\% -- 89.999\% \\
        B  & (3.0) & 84\% -- 86.999\% (Good/Better than Average) \\
        B- & (2.7) & 80\% -- 83.999\% \\
        C+ & (2.3) & 77\% -- 79.999\% \\
        C  & (2.0) & 74\% -- 76.999\% (Competent/Average) \\
        C- & (1.7) & 70\% -- 73.999\% (Minimum Passing for Business Core) \\
        D+ & (1.3) & 67\% -- 69.999\% \\
        D  & (1.0) & 64\% -- 66.999\% \\
        D- & (0.7) & 60\% -- 63.999\% \\
        F  & (0.0) & 0\% -- 59.999\% (Failing) \\
        \bottomrule
    \end{tabular}
\end{center}
\textit{Note: Grading policies of the CU Denver Business School state that the average GPA across all students in an undergraduate class should generally fall within the range of 2.3 (C+) to 3.0 (B). If necessary, grading scales may be slightly adjusted formulaically to align within this institutional guideline.}

\section*{Course Schedule \& Modular Breakdown}
The course is structured across 16 weeks, broken down into 8 core instructional content modules:

\begin{itemize}[leftmargin=*]
    \item \textbf{Module 1: Creating a Worksheet and a Chart (Weeks 1 \& 2)}
    \begin{itemize}
        \item \textit{Key Topics:} Microsoft Office ribbon components, Excel worksheet structure, simple functions/text/numbers, cell styling, and Pie Charts.
        \item \textit{Project \& Assignments:} Real Estate budget and chart project; Apply Your Knowledge (Changing Values); Extend Your Knowledge (Styles and Formatting); Expand Your World (Loan Calculator); In the Lab (Comparing Laptops).
    \end{itemize}
    
    \item \textbf{Module 2: Formulas, Functions, and Formatting (Weeks 3 \& 4)}
    \begin{itemize}
        \item \textit{Key Topics:} Flash Fill, math formulas, workbook themes, Conditional Formatting, and printing options.
        \item \textit{Project \& Assignments:} Klapore Engineering Salary Report; Apply Your Knowledge (Cost Analysis); Extend Your Knowledge (Kalto Security Outlet tracking); Expand Your World (4-Year College Cost Calculator); In the Lab (Cell Phone Service Summary).
    \end{itemize}

    \item \textbf{Module 3: Working with Large Worksheets, Charting, and What-If Analysis (Weeks 5 \& 6)}
    \begin{itemize}
        \item \textit{Key Topics:} Text rotations, advanced formatting, Absolute vs. Relative Cell References, logical tests (IF functions), Sparkline charts, Format Painter, What-If Analysis, freezing/unfreezing panes.
        \item \textit{Project \& Assignments:} Manola Department Stores Six-Month Financial Projection; Apply Your Knowledge (Logical Tests and Absolute Referencing); Extend Your Knowledge (Fill Handle and Nested IF Statements); Expand Your World (Census Data); Quiz.
    \end{itemize}

    \item \textbf{Module 4: Financial Functions, Data Tables, and Amortization (Weeks 7 \& 8)}
    \begin{itemize}
        \item \textit{Key Topics:} Financial Functions (PMT, etc.), Data Tables, Amortization Schedules, PV vs. FV analysis, protecting/unprotecting worksheets, hiding/unhiding sheets.
        \item \textit{Project \& Assignments:} Cranfod Credit Union Mortgage Payment Calculator; Apply Your Knowledge (Loan Payments); Extend Your Knowledge (Retirement Planning); Expand Your World (Down Payment Options); Critical Thinking Questions.
    \end{itemize}

    \item \textbf{Module 5: Working with Multiple Worksheets and Workbooks (Weeks 9 \& 10)}
    \begin{itemize}
        \item \textit{Key Topics:} Formatting and consolidation techniques, Linear Series fills, custom format codes, custom cell styles, 3-D cell references, and 3-D Pie Charts.
        \item \textit{Project \& Assignments:} M\&S Provisions project; Apply Your Knowledge (Consolidating Payroll); Extend Your Knowledge (Custom Format Codes); Expand Your World (Weather Data tracking).
    \end{itemize}

    \item \textbf{Module 6: Creating, Sorting, and Querying a Table (Weeks 11 \& 12) --- MID-TERM EXAM}
    \begin{itemize}
        \item \textit{Key Topics:} Table creation/manipulation, calculated columns, XLOOKUP function, sorting, and querying/searching text via AutoFilter.
        \item \textit{Project \& Assignments:} Rating Bank Account Managers project; Apply Your Knowledge (Table with Conditional Formatting); Extend Your Knowledge (Advanced Functions); Expand Your World (Converting Files); Mid-Term Exam.
    \end{itemize}

    \item \textbf{Module 7: Creating Templates, Importing Data, Pivot Tables, and Macros (Weeks 13 \& 14)}
    \begin{itemize}
        \item \textit{Key Topics:} Excel Templates, SmartArt graphics, image insertion, Pivot Tables, Pivot Charts, and Macro basics.
        \item \textit{Project \& Assignments:} Mayor Insurance project; Apply Your Knowledge (Template Consolidated Workbook); Extend Your Knowledge (SmartArt Org Chart and Image); Expand Your World (PivotTable and PivotChart creation).
    \end{itemize}

    \item \textbf{Module 8: Data Analysis, Charting, and Dashboards using Tableau (Weeks 15 \& 16) --- FINAL EXAM}
    \begin{itemize}
        \item \textit{Key Topics:} Data management concepts, data visualization basics, Tableau data loading and handling, building interactive dashboards, and storytelling with data.
        \item \textit{Project \& Assignments:} Apply Your Knowledge (Bar Charts in Tableau); Extend Your Knowledge (Data Analysis and Interpretation in Tableau); Course Review; Final Exam.
    \end{itemize}
\end{itemize}

\section*{Strict Course Policies}

\subsection*{Late Work Policy}
All homework files must be submitted through the appropriate Canvas module page. Assignments are due by 11:59 PM MST on the specified dates. Assignments posted after the indicated due dates will be subject to an immediate loss of \textbf{10\% of available points for each day late}. No assignment will be accepted for a grade after 9 days late. On the 10th day late, a score of \textbf{0} is finalized. No assignments can be accepted after 11:59 PM on the final day of the semester class session.

\subsection*{Academic Honesty \& Security Statement}
Students are bound to comply with the CU Denver Student Code of Conduct against plagiarism, cheating, fabrication, and falsification. The incorporation of another's work into your submission without explicitly authorized parameters is academic dishonesty.
\\
\textbf{Mandatory Syllabus Declaration:} When you log in to your courses using your CU Denver login credentials, you are certifying that the person logging in is you. As such, providing your login credentials to any other individual and allowing them to log in to a course on your behalf is considered academic dishonesty. Sharing files, copying formula text, or outsourcing graded dashboard setups will result in an automatic grade of "F" in the course and a direct escalation to the Student Conduct Committee.

\subsection*{Best Mode of Communication}
The best mode of communication is the \textbf{messaging function directly within Canvas}. Canvas messages send an integrated copy to the instructor's inbox clearly marked with the course number. Do \textbf{not} email files directly from your personal or student email programs; Outlook automated protocols regularly block attachments containing macro-embedded spreadsheets or code files, meaning your communication could be completely missed.

\subsection*{Attendance, Laptop, and Mobile Device Rules}
Participation heavily influences mastery of technical tools. Attendance will be tracked. If you are ill, you must email the professor via Canvas in advance.
\\
Laptops are authorized strictly for active note-taking and working on live in-class modules. Laptops shall \textbf{not} be used for social media, internet surfing, or messaging during class. \textbf{Mobile phones and texting are strictly prohibited.} Students who use their cell phones during active class sessions will be asked to leave immediately and will receive an unexcused absence for that calendar day.

\subsection*{Visible ID Badge Requirement}
Safety continues to be a priority focus for the CU Denver campus community. All students, faculty, and staff in the Business School building are required to wear their official University ID badge in a visible location (such as on a lanyard, card holder, or clothing clip) at all times while inside the building.

\subsection*{Grades of Incomplete}
Incomplete ("I") grades are granted strictly at the instructor's final authority in unexpected emergencies. An incomplete is considered only if the student has successfully completed a minimum of \textbf{75\% of all course coursework} with passing grades, and the remaining portion can be independently wrapped up during the subsequent semester.

\section*{Student Success Coaching (SSC) Program}
This course features an integrated Student Success Coach, \textbf{Abdelkarim Mansour}, dedicated entirely to supporting student performance and navigation through complex tools. 
\\
The Success Coach will actively reach out if you drop an active class attendance day, fail to turn in a scheduled file, or if your classroom average in Canvas falls below a "C" grade. 
\\
\textbf{The Learning Experience Rule:} If you receive a "C" or below on any major assignment (representing 5\% or more of the course grade), you are \textbf{required} to meet with the Success Coach (virtually or in person) \textbf{BEFORE the next major assignment is due}. During this interactive session, the coach will assist you with understanding the mistakes made. Once you successfully complete this review meeting, you will receive \textbf{5\% of the points back} on that poor assignment score.

\section*{University Resources \& Access Accommodation}
\begin{itemize}[leftmargin=*]
    \item \textbf{Disability Access (DRS):} Students entitled to academic adjustments must register directly with Disability Resources and Services in Student Commons Room 2116 (303-315-3510). Approved accommodation letters must be shared directly with the instructor to execute support.
    \item \textbf{Learning Resources Center (LRC):} Offers 24/7 online tutoring services covering Excel/Tableau applications, plus in-person time management and study workshops (Student Commons \#1231, \url{https://www.ucdenver.edu/lrc}).
    \item \textbf{Mental Health \& Support:} The Student and Community Counseling Center (Tivoli Student Union \#454, 303-315-7270) provides cost-free, completely confidential psychological and mental health counseling.
    \item \textbf{Campus Safety Concerns:} If you have an immediate concern about the behavior or safety of anyone, submit an anonymous alert directly to the Campus Assessment, Response \& Evaluation (CARE) Team at 303-315-7306.
\end{itemize}

\end{document}
